\section{Detailed derivations and convention checks}
\label{app:derivations}

\subsection{DFPT response factorization in a scalar system}

Starting from \cref{eq:vksresponse,eq:chisresponse},
\begin{align}
  \delta V_{\ks}
  &=\delta V_{\mathrm{ion}}
   +(v+f_{\mathrm{xc}})\chi_s\delta V_{\ks},\\
  \bigl[1-(v+f_{\mathrm{xc}})\chi_s\bigr]\delta V_{\ks}
  &=\delta V_{\mathrm{ion}}.
\end{align}
With $P=\chi_s/(1-f_{\mathrm{xc}}\chi_s)$,
\begin{align}
  (1-vP)(1-f_{\mathrm{xc}}\chi_s)
  &=1-f_{\mathrm{xc}}\chi_s
   -v\frac{\chi_s}{1-f_{\mathrm{xc}}\chi_s}
       (1-f_{\mathrm{xc}}\chi_s)\\
  &=1-(v+f_{\mathrm{xc}})\chi_s.
\end{align}
Thus
\begin{equation}
  g^{\dfpt}
  =\frac{g^0}{(1-vP)(1-f_{\mathrm{xc}}\chi_s)}.
\end{equation}
Comparing with $g^{\mathrm{phys}}=g^0z\Gamma_\rho/(1-vP)$ yields
$z\Gamma_\rho=P/\chi_s$.  This algebra establishes the matching condition under
the scalar definitions; it does not establish that the condition is exact.

\subsection{Uniform-potential cancellation at a quasiparticle pole}

Let the pole satisfy
\begin{equation}
  E_u-\varepsilon-Du-\Sigma_0(E_u-Du)=0.
\end{equation}
Differentiate at $u=0$:
\begin{equation}
  \frac{dE}{du}-D
  -\partial_\omega\Sigma_0\left(\frac{dE}{du}-D\right)=0.
\end{equation}
Provided the pole is well-defined,
$(1-\partial_\omega\Sigma_0)(dE/du-D)=0$, so $dE/du=D$.  In the alternative
form where $\partial_u\Sigma=-D\partial_\omega\Sigma$, the same result is
$z(D+\partial_u\Sigma)=D$ with
$z=(1-\partial_\omega\Sigma)^{-1}$.  The result follows from the common shift of
frequency and chemical potential, not from weak interaction.

\subsection{Limit ordering in the Ward identity}

Expand the right-hand side of \cref{eq:ward} for small four-momentum:
\begin{equation}
  G^{-1}(k+q)-G^{-1}(k)
  =\ii\nu\,\partial_{\ii\omega}G^{-1}
   +\vq\cdot\nabla_{\vk}G^{-1}+\order(q^2).
\end{equation}
At $\vq=0$, matching the coefficient of $\ii\nu$ gives the scalar dynamic
vertex.  At $\nu=0$, matching the coefficient of $\vq$ gives the longitudinal
current vertex.  A static scalar perturbation also changes the distribution and
self-consistent field, producing the compressibility vertex.  This explicit
expansion shows why swapping the two limits changes the conclusion.

\subsection{Matrix generalization}

In a crystal, $v$, $P$, $\chi_s$, $f_{\mathrm{xc}}$, and $\epsilon$ are matrices
in reciprocal vectors $\bm G,\bm G'$.  Products do not generally commute, so the
scalar factorization in \cref{eq:factorization} cannot be transferred by simply
replacing every scalar with a matrix.  A crystal implementation must retain
matrix ordering and local-field components, or derive a controlled projected
version of the relation.  This is one of the first theoretical tasks of the
project.
